\section{Search and Retrieval}
\label{sec:app_retrieval}

\subsection{API Contract}

Search is exposed via a single endpoint, \texttt{POST /search}, which accepts a JSON body with three optional fields:

\begin{itemize}
    \item \texttt{query} (string, default \texttt{""}): the text query, in any of the 19 supported languages.
    \item \texttt{image\_base64} (string, optional): a base64-encoded image for CLIP visual search.
    \item \texttt{top\_k} (integer, default 20): number of results to return.
\end{itemize}

The response is a JSON object containing a \texttt{results} array of hydrated item records (ASIN, title, author, genre, description, cover URL, relevance score, and per-modality similarity scores when applicable), a \texttt{total} count, a \texttt{live\_encoding} flag, and the echoed \texttt{query} string.
An optional \texttt{debug=true} query parameter appends a \texttt{\_debug\_timings} dictionary to the response, reporting per-stage latency in milliseconds (translation, text encoding, CLIP encoding, FAISS search, BM25 re-ranking, metadata hydration, and total).

\subsection{Retrieval Pipeline}

The internal pipeline executed by the \texttt{POST /search} handler is:

\begin{enumerate}
    \item \textbf{Language detection and translation}: the \texttt{lingua} detector identifies the input language in under 1\,ms. Non-English queries are translated to English by the NLLB-200 pipeline (Section~\ref{sec:nllb_pipeline}).
    \item \textbf{Parallel encoding}: text and image are encoded concurrently via an \texttt{anyio} task group. Each encoder runs in a background thread to avoid blocking the async event loop.
    \item \textbf{FAISS HNSW search}: the text vector is queried against \nolinkurl{bge_index_hnsw.faiss}; the image vector (if present) is queried against the CLIP HNSW index. Both retrieve 100 candidates.
    \item \textbf{Tantivy BM25 re-ranking}: the original query string is searched against the Tantivy full-text index to retrieve a ranked keyword hit list.
    \item \textbf{RRF fusion}: the HNSW and BM25 candidate lists are fused by Reciprocal Rank Fusion~\cite{cormack2009rrf} to produce a single ranked list.
    \item \textbf{Metadata hydration}: each ASIN in the final list is resolved against the in-memory DataFrame to attach display metadata.
\end{enumerate}

\subsection{Profile-Based Retrieval}

The \texttt{GET /recommend} endpoint uses a separate retrieval path.
For authenticated warm users, the profile manager fetches the user's most recent click sequence and computes a mean-pool BGE-M3 profile vector.
Pipeline~A queries the Cleora FAISS index (\texttt{IndexFlatIP}) with this profile vector, returning the 50 nearest co-purchase neighbours.
Pipeline~B retrieves item embeddings for the candidate set from the flat BGE-M3 index using the \texttt{reconstruct(i)} operation (HNSW indices do not support reconstruction, which is why the flat index is maintained separately), and applies the content veto filter before returning the DIF-SASRec top items.
